\documentclass[10pt,twocolumn]{article}
\usepackage[margin=0.75in,columnsep=0.25in]{geometry}
\usepackage{mathptmx}
\usepackage{amsmath}
\usepackage{graphicx}
\usepackage{booktabs}
\usepackage{makecell}
\usepackage{xcolor}
\usepackage{amssymb}
\usepackage[hidelinks]{hyperref}

% Tables and figures come straight from eval/analyze.py:
%   eval/.venv/bin/python eval/analyze.py runs/<session> --out results/final
\newcommand{\res}{../results/final}
% Result-dependent statements still to be written from the final data.
\newcommand{\todo}[1]{\textcolor{red}{\textbf{TODO:} #1}}

\title{\textbf{LoopShield: Layered Sensor Spoofing Detection\\for Real-Time Controllers}}
\author{Saad Kabir Uddin}
\date{}

\begin{document}
\maketitle

\begin{abstract}
Networked controllers increasingly receive their sensor readings over
general-purpose links that an attacker on the network can tamper with. We ask
whether a resource-constrained real-time controller can detect sensor spoofing
on that link---at both the protocol layer and the physical layer---without
missing its control deadlines. We build a closed-loop testbed in which an
ESP32 running FreeRTOS regulates a simulated room's temperature over MQTT,
an on-path attacker mounts five attacks (bias injection, slow ramp, replay,
forgery and flooding), and two independent, separately switchable defense
layers guard the controller: per-message authentication and freshness checks,
and a physics-based detector combining a model invariant with CUSUM.
Across the full attack $\times$ defense matrix we measure detection rate, false
positives, time to detect, worst true-temperature deviation and the real-time
cost of each layer. \todo{one-sentence headline result from the matrix and timing table.}
\end{abstract}

\section{Introduction}
Cyber-physical systems close a control loop through a network: a sensor
reports, a controller computes, an actuator acts. If the sensor's messages can
be altered in transit, the controller can be steered into driving the physical
process anywhere the attacker likes while believing everything is
fine---the core of the Stuxnet attack, which replayed normal-looking readings
while the plant was being damaged.

Two families of defenses exist. \emph{Protocol-layer} defenses (message
authentication codes, sequence numbers, freshness windows) stop an attacker who
lacks the key, at the price of cryptography on every message. \emph{Physics-layer}
defenses check that measurements are consistent with what the physical process
can do~\cite{cardenas2011,giraldo2018}; they need no key and catch a compromised
sender, but they are statistical, so they trade detection delay against false
alarms. Both cost CPU time, and a real-time controller has a fixed budget: a
defense that makes the control task miss its deadlines is itself a
denial-of-service vector.

This paper asks: \emph{can a resource-constrained real-time controller run both
kinds of defense and still meet its deadlines, and which layer catches which
attack?} Our contributions are:
\begin{itemize}
  \item an open testbed (ESP32/FreeRTOS controller in Wokwi, C\# plant
        simulator and on-path attacker, automated harness) in which every
        defense can be toggled at runtime so that each layer's contribution is
        isolated;
  \item a lightweight physics detector (model observer, slope invariant,
        two-sided CUSUM) that fits in its own RTOS task;
  \item the attack $\times$ defense matrix and per-layer timing, including under
        flooding.
\end{itemize}

\section{System and Threat Model}
\paragraph{Plant.} A single room with a heater, modelled as
$\dot T = \big((T_{amb} - T) + K u\big)/\tau$ with $T_{amb}=15\,^\circ$C,
$K=20\,^\circ$C, $\tau=60$\,s and heater duty $u\in[0,1]$. The plant runs as a
separate program, publishes the temperature plus Gaussian noise
($\sigma=0.05\,^\circ$C) every 100\,ms, and logs the \emph{true} temperature,
which gives every experiment ground truth the attacker cannot touch.

\paragraph{Controller.} An ESP32 (Xtensa LX6, 160\,MHz, dual core) running
ESP-IDF~v6.1/FreeRTOS, simulated in Wokwi. A PI controller
($K_p=0.5$, $T_i=60$\,s, conditional-integration anti-windup) runs as a periodic
task every 100\,ms with an implicit deadline equal to its period, regulating the
room to $22\,^\circ$C.

\paragraph{Link.} All traffic is MQTT over TCP. The controller connects to the
broker through a proxy, which is where the attacker sits.

\paragraph{Attacker.} An on-path man-in-the-middle on the controller's link
(as obtained by ARP spoofing or a rogue access point on an unencrypted network).
It can read, drop, modify, delay, replay and inject any message, but it
compromises neither node's firmware and does not know the shared key. It
implements five attacks on sensor messages:
\emph{bias} adds $-3\,^\circ$C; \emph{slow ramp} adds an offset growing at
$-0.01\,^\circ$C/s; \emph{replay} records 30\,s of genuine readings and plays
them back in place of live ones while the room is genuinely cooled by an
ambient drop to $5\,^\circ$C (the Stuxnet pattern); \emph{forgery} injects,
after every genuine reading, a fabricated one claiming $18\,^\circ$C with the
next sequence number and a random MAC; \emph{flooding} injects 1000 random
sensor messages per second. Out of scope: physical attacks on the sensor,
compromised endpoints, and attacks on the out-of-band configuration channel
used by the test harness.

\section{Defenses}
Every sensor message is a fixed 28-byte frame: version, type, flags, a 32-bit
sequence number, a 64-bit sender timestamp, the float reading and an 8-byte
truncated HMAC-SHA256. All fields are present in every configuration, so
switching a defense on changes only the \emph{checking}, never message size or
parsing cost.

\subsection{Protocol layer}
Three checks, each switchable, run in the MQTT receive path, in this order:
\begin{enumerate}
  \item \textbf{MAC}: HMAC-SHA256 truncated to 64 bits (PSA Crypto), compared in
        constant time. It runs first so that unauthenticated messages can never
        advance the sequence or time state.
  \item \textbf{Sequence}: strictly increasing sequence numbers.
  \item \textbf{Freshness}: sender and receiver clocks are not synchronised (and
        the simulator runs a few percent slower than real time). The offset
        $o = t_{local} - t_{sender}$ is network delay plus clock offset, and
        delaying a message can only \emph{increase} it. The receiver therefore
        tracks the smallest offset seen, $o_{min}$, and rejects a message if
        $o > o_{min} + 250\,\text{ms} + \rho\,\Delta t$, where $\Delta t$ is the
        time since $o_{min}$ was observed and $\rho=1\%$ absorbs a receiver clock
        that runs fast. Sender timestamps must also strictly increase, which
        rejects replays and duplicates. Section~\ref{sec:search} shows why the
        reference must be the minimum and not the last message.
\end{enumerate}
If no message is accepted for five periods, the controller declares the sensor
lost and holds its output rather than integrating on stale data.

\subsection{Physics layer}
A separate RTOS task receives each period's accepted measurement $y_k$ and the
heater command $u_{k-1}$ that acted during the period.
\paragraph{Observer.} $\hat T$ follows the nominal model driven by $u$, with a
small correction $\ell=0.01$ toward $y$. The innovation
$r_k = y_k - \hat T_{k|k-1}$ is zero-mean noise when the sensor is honest. A bias
shifts it immediately; a slow ramp produces a persistent offset, because the
closed loop drives $y$ to the setpoint while $u$, which the observer sees,
follows the \emph{true} temperature.
\paragraph{CUSUM.} Two-sided: $S^\pm_k = \max(0, S^\pm_{k-1} \pm r_k - k)$,
alarm when $S^+$ or $S^-$ exceeds $h$. We use $k=0.05$ and $h=4$, chosen in a
host simulation (zero false alarms in 10 simulated hours, a
$0.005\,^\circ$C/s ramp detected in about 50\,s); Fig.~\ref{fig:roc} sweeps $h$
on the testbed data.
\paragraph{Invariant.} The least-squares slope of $y$ over the last 2\,s must
lie in the range the model allows for the observed $u$ and \emph{any} ambient
temperature in $[5, 25]\,^\circ$C, plus a 3$\sigma$ tolerance. ``Heater on but
temperature flat'' and ``faster change than any heater can produce'' both
violate it; three consecutive violations raise an alarm.
\paragraph{Mitigation.} On a physics alarm the observer stops trusting $y$ and
the controller regulates the model estimate $\hat T$ instead.

\section{Implementation}
The firmware separates concerns by task so that timing measurements mean what
they say (Table~\ref{tab:tasks}). The control task performs no I/O: it reads the
latest accepted sample from a length-one queue and hands the actuator command
to a low-priority transmit task. Protocol checks run in the MQTT task on the
network core, so their cost under flooding appears as interference rather than
as control-task execution time. Wi-Fi power saving is disabled: its default
modem sleep wakes only at DTIM beacons (a 300--400\,ms listen interval in our
setup), which can delay received readings by several control periods.

\begin{table}[t]
\centering\small
\caption{Firmware task layout.}
\label{tab:tasks}
\begin{tabular}{llcc}
\toprule
Task & Work & Core & Prio \\
\midrule
control & PI / safe mode, 100\,ms & 1 & 10 \\
detector & observer, CUSUM, invariant & 1 & 8 \\
MQTT & receive + protocol checks & 0 & 5 \\
tx & publish actuator, telemetry & 0 & 4 \\
Wi-Fi / lwIP & network stack & 0 & 23 / 18 \\
\bottomrule
\end{tabular}
\end{table}

Each control job $j$ is released at $t_0 + jP$; its response time is measured
from that nominal release to completion, so delayed wake-ups count. A deadline
miss is a response time above $P$. Every period the controller publishes a
telemetry record (measurement, command, mode, observer state, CUSUM statistics,
cumulative accept/reject counters, control and detector execution and response
times, and per-period protocol-check cost); a summary with cumulative worst
cases is published every 10\,s.

The protocol checks, detector and PI law are portable C with no ESP-IDF
dependencies. They are unit-tested on the host, and a host closed-loop
simulation runs the same code faster than real time for tuning.

\section{Evaluation Methodology}
A harness runs each (attack, configuration, trial) cell automatically: it
starts the attacker (as a transparent relay in attack-free runs, so every run
has the same network path), sends the configuration to the controller and
waits for its acknowledgement (which also resets PI, detector, protocol state
and statistics), and starts the plant warm at the setpoint. Each run is a
30\,s warm-up, a 120\,s attack window and 30\,s of recovery. Configurations are
\emph{none}, \emph{protocol only}, \emph{physics only} and \emph{both}, plus a
benign \emph{disturbance} scenario (a real ambient drop with no attacker) to
measure false positives on a legitimate change. \todo{number of trials.}

\paragraph{Metrics.} \emph{Detection}: an enabled layer flags the attack inside
the attack window (first rejected message or first physics alarm).
\emph{Time to detect}: from attack start to that flag. \emph{False positive}:
any flag before the attack (or anywhere in attack-free runs). \emph{Max
deviation}: worst $|T_{true}-22\,^\circ\text{C}|$ from attack start to run end,
from the plant's ground truth. \emph{Timing}: control execution and response
times, detector execution, protocol-check cost per message, and deadline misses.

\section{Results}

\begin{table*}[t]
\centering\small
\caption{Attack $\times$ defense matrix. Each cell: detection rate / mean time to
detect, then the worst true-temperature deviation from the setpoint (and the
false-positive rate when non-zero). Attack-free rows show false-positive rate and
deviation. $^*$Disturbance is a real ambient drop, not an attack.}
\label{tab:matrix}
\input{\res/matrix.tex}
\end{table*}

\begin{figure*}[t]
\centering
\includegraphics[width=\textwidth]{\res/temp_grid.png}
\caption{True room temperature under each attack and defense configuration
(trial 1). Shaded: attack window. Green band: $\pm 0.5\,^\circ$C.}
\label{fig:temps}
\end{figure*}

\paragraph{Which layer catches which attack.} Table~\ref{tab:matrix} and
Fig.~\ref{fig:temps}. \todo{describe from the data. Expected from design and the
host simulation, to be confirmed or refuted: the protocol layer rejects bias,
ramp and forgery by MAC and replay by sequence/freshness; physics catches bias,
forgery and flooding within a sample and the ramp after it accumulates, but not
replay, because replayed readings are physically consistent; the disturbance
raises a physics false positive, because the model cannot distinguish an
unknown ambient change from spoofing.}

\subsection{Protocol state-machine search}\label{sec:search}
We model the receiver's checks as a state machine and search exhaustively
(memoised depth-first search, 50\,ms steps over 1.4\,s, after three on-time
messages) over attacker action sequences: deliver nothing; deliver any genuine
message emitted so far, early, late, again or out of order; or forge a message
with the next or a far-ahead sequence number. A delivery the firmware's own
check function accepts is malicious if it is forged, a replay (not newer than
an accepted message) or stale (more than 500\,ms old on arrival).
Table~\ref{tab:search} summarises the result for all eight check combinations.

\begin{table}[t]
\centering\small
\caption{Shortest malicious sequence accepted, per check combination
(\checkmark{} = none found within the horizon).}
\label{tab:search}
\begin{tabular}{lccc}
\toprule
Checks & Forged & Replay & Stale \\
\midrule
none            & 1 step & 1 step & 1 step \\
seq             & 1 step & \checkmark & 1 step \\
fresh, seq+fresh& 1 step & \checkmark & \checkmark \\
mac             & \checkmark & 1 step & 1 step \\
mac+seq         & \checkmark & \checkmark & 1 step \\
mac+fresh, all  & \checkmark & \checkmark & \checkmark \\
\bottomrule
\end{tabular}
\end{table}

Three findings. (i) Our first freshness check compared each message only with
the \emph{last accepted} one. The search found an \emph{incremental delay}
attack against it: delay message $n$ by 250\,ms, $n{+}1$ by 500\,ms, $n{+}2$ by
550\,ms. Each step looks fresh relative to its predecessor, so the delay grows
without bound. Anchoring to the minimum offset caps the added delay at
$250\,\text{ms} + \rho\,\Delta t$, and the search then finds nothing. (ii) Once
freshness enforces strictly increasing authenticated timestamps, the sequence
check is redundant in this model. (iii) Every combination trusts the
\emph{first} message after a reset, however old. A cold-start receiver accepts
an hour-old genuine message, so a deployment needs an initial challenge or
synchronised time.

\paragraph{Seq numbers without a MAC are not enough.} With sequence checking but
no MAC, a forged message carrying the next sequence number is accepted and
then causes the genuine message with that number to be rejected, locking the
real sensor out (a unit test in the repository shows this). \todo{whether the
testbed shows this in the protocol-only forgery cell.}

\begin{table}[t]
\centering\scriptsize
\caption{Real-time cost per configuration, without attack and under flooding.
Control execution covers one job; WCRT is the worst control response time over
the whole run (deadline 100\,ms).}
\label{tab:timing}
\resizebox{\columnwidth}{!}{\input{\res/timing.tex}}
\end{table}

\paragraph{Real-time cost.} Table~\ref{tab:timing} and Fig.~\ref{fig:timing}.
\todo{HMAC cost per message; detector execution per period as a fraction of the
100\,ms budget; WCRT and deadline misses with and without defenses; flooding:
whether per-message checking on the network core interferes with the control
task on the other core, and the protocol-check CPU share at 1000\,msg/s.}

\begin{figure}[t]
\centering
\includegraphics[width=\columnwidth]{\res/timing.png}
\caption{Execution cost of each layer and worst-case control response time.}
\label{fig:timing}
\end{figure}

\begin{figure}[t]
\centering
\includegraphics[width=\columnwidth]{\res/roc.png}
\caption{CUSUM ROC over the threshold $h$, recomputed offline from the logged
measurements and commands (run level; negatives are attack-free segments).}
\label{fig:roc}
\end{figure}

\paragraph{Detection threshold.} Fig.~\ref{fig:roc}. \todo{shape of the ramp ROC
and whether $h=4$ sits at the knee.}

\section{Discussion and Limitations}
\textbf{The layers are complementary.} The protocol layer is exact against an
outsider without the key but blind to a compromised sender; the physics layer
needs no key but is blind to physically plausible lies such as replay, and it
cannot tell spoofing from a real change its model does not know (the
disturbance row). A constant bias is, in steady state, only detectable at all
because the model assumes a known ambient temperature.
\textbf{Simulation.} Wokwi is not cycle-accurate and ran about 2\% slower than
real time; absolute execution times are indicative, and relative comparisons
between configurations are the meaningful result. \todo{confirm with the
measured spread.} \textbf{Keys.} A single pre-shared lab key; provisioning and
rotation are out of scope. \textbf{Configuration channel.} The harness
reconfigures the controller over an unauthenticated topic, which the attacker
in our model does not touch; a deployment would authenticate it or remove it.

\section{Related Work}
Physics-based attack detection in CPS is surveyed by Giraldo et
al.~\cite{giraldo2018}. C\'ardenas et al.~\cite{cardenas2011} apply CUSUM to
process-control attacks, and Urbina et al.~\cite{urbina2016} show that stateful
detectors such as CUSUM limit the impact of stealthy attacks better than
stateless ones. Mo and Sinopoli~\cite{mo2009} analyse replay attacks and
propose watermarking the control input. Truncated MACs with freshness values
are standard in automotive networks (AUTOSAR SecOC). We combine both layers on
one constrained real-time controller and measure their cost against its
deadlines.

\section{Conclusion}
\todo{answer the research question in two or three sentences from the results.}

\paragraph{Reproducibility.} Firmware, plant, attacker, harness, analysis and
raw run logs: \url{https://github.com/saadkabiruddin/spoofguard}.

\begin{thebibliography}{9}\small
\bibitem{cardenas2011} A.~A. C\'ardenas, S.~Amin, Z.-S. Lin, Y.-L. Huang, C.-Y. Huang,
  and S.~Sastry, ``Attacks against process control systems: risk assessment,
  detection, and response,'' in \emph{Proc. ACM ASIACCS}, 2011.
\bibitem{giraldo2018} J.~Giraldo, D.~Urbina, A.~Cardenas, J.~Valente, M.~Faisal,
  J.~Ruths, N.~O. Tippenhauer, H.~Sandberg, and R.~Candell, ``A survey of
  physics-based attack detection in cyber-physical systems,'' \emph{ACM
  Computing Surveys}, vol.~51, no.~4, 2018.
\bibitem{urbina2016} D.~I. Urbina, J.~A. Giraldo, A.~A. Cardenas, N.~O. Tippenhauer,
  J.~Valente, M.~Faisal, J.~Ruths, R.~Candell, and H.~Sandberg, ``Limiting the
  impact of stealthy attacks on industrial control systems,'' in \emph{Proc. ACM
  CCS}, 2016.
\bibitem{mo2009} Y.~Mo and B.~Sinopoli, ``Secure control against replay attacks,''
  in \emph{Proc. Allerton Conference on Communication, Control, and Computing}, 2009.
\end{thebibliography}

\end{document}
